%%
% 第一章 绪论
% 面向评委介绍研究背景、国内外现状、P450 家族上的特殊挑战、本研究目标与组织
%%

\chapter{绪论}
\label{cha:introduction}

\section{研究背景与意义}
\label{sec:background}

酶的底物特异性是酶功能的核心属性，指酶识别并选择性作用于特定底物的能力。这种特异性根植于酶活性位点的三维结构与反应过渡态之间的精细耦合：底物进入活性位点时，需在空间取向、电荷分布与氢键网络等多个维度与酶精确匹配，方能触发高效催化。迄今为止已确定序列的酶达数百万之多，其中绝大多数缺乏可靠的底物特异性注释，严重制约了生物催化、合成生物学与药物研发等应用领域中酶工具的挑选与设计。

在众多酶家族中，细胞色素 P450 超家族（Cytochrome P450，CYP）因其在生命活动中广泛而关键的代谢功能而格外受关注。P450 参与约 75\% 临床上市药物的第 I 相代谢过程\cite{Nelson2009}，是药物动力学、药物--药物相互作用（drug--drug interaction，DDI）与个体代谢差异的核心决定因素之一：CYP3A4 参与辛伐他汀、阿托伐他汀、咪达唑仑等多类常用药物的代谢；CYP2D6 的遗传多态性则直接影响可待因、他莫昔芬等前药的活化效率。在植物界，P450 在紫杉醇、青蒿素、人参皂苷等重要天然药物的生物合成路径中承担多步氧化修饰\cite{Nelson2004}。在农业与环境层面，P450 介导的氧化修饰决定了许多农药与外源化合物的降解命运。因此，精准预测特定 P450 酶与小分子底物之间的特异性，对药物代谢筛选、合成生物学改造与农业化学品管理均具有直接的应用价值。

\section{国内外研究现状评述}
\label{sec:related_work}

传统的酶底物特异性研究依赖昂贵的体外酶学实验（如 $K_m$、$k_{cat}$ 测定）与底物筛选，周期长且难以规模化。随着蛋白质序列与结构数据规模的指数增长，计算预测方法逐步兴起。近年来典型的方法沿两条主线发展。一类方法基于一维蛋白序列与蛋白语言模型：CLEAN 基于蛋白语言模型 ESM-2\cite{ESM2023} 结合对比学习预测酶的 EC 编号；ESP 则在此基础上引入分子指纹通道，直接预测酶-底物结合活性。另一类方法基于分子图与蛋白-化合物对建模（compound--protein interaction，CPI），通过图神经网络分别编码蛋白侧与底物侧，再以多种交互机制融合预测。这两类方法主要依赖一维序列或二维分子图表示，难以捕捉底物进入活性位点过程的三维本质，对同一功能类内不同酶的底物偏好区分能力有限。

2025 年发表于 \textit{Nature} 的 EZSpecificity 模型\cite{EZSpec2025} 代表了当前 SOTA 水平。该方法的创新体现于数据与模型两个层面。数据层面，作者构建了首个结构级酶--底物复合物数据库 ESIBank，整合了 BRENDA\cite{Chang2021} 与 UniProt\cite{UniProt2025} 的序列与功能标注数据、AlphaFold 数据库\cite{Varadi2024} 与 AlphaFill\cite{AlphaFill2023} 的预测结构与辅因子回填，以及经 AutoDock-GPU\cite{Vina2010,Vina2021} 对接得到的总共 323{,}783 对酶--底物复合物结构。模型层面，EZSpecificity 采用 SE(3)-等变图神经网络编码活性位点三维微环境，以 ESM-2\cite{ESM2023} 提供残基级序列表征，并通过双层交叉注意力机制建模酶--底物原子级交互。在 8 种卤化酶 $\times$ 78 种底物的湿实验验证上，EZSpecificity 的 Top-1 准确率达到 91.7\%，远高于此前 SOTA 方法 ESP 的 58.3\%\cite{EZSpec2025}。论文声称该方法具有通用性，可跨多种蛋白家族预测底物特异性。

\section{通用预测模型在 P450 上的局限与待解决问题}
\label{sec:limitations}

对 EZSpecificity 通用性声明的一个自然追问是：对于训练集中覆盖稀缺但实际应用极重要的酶家族，该声明是否依然成立？EZSpecificity 论文重点验证了酯酶、糖基转移酶、硝酸酶、磷酸酶、硫解酶、DUF 蛋白六个代表性家族及卤化酶，这些家族在 ESIBank 训练集中均具有充分的数据覆盖。

细胞色素 P450 恰好是检验上述通用性声明的理想试金石。从应用价值看，P450 在药物代谢与植物次生代谢中的核心地位如前所述；从结构特点看，P450 以血红素铁（heme-Fe）为催化中心，具有深长的底物通道（substrate access channel）与开合构象变化，与 ESIBank 中其他家族的活性位点形态显著不同；然而从数据覆盖看，对 ESIBank 中全部 25{,}225 个酶逐条进行 InterPro\cite{Blum2025} 功能域验证后，真正的 P450 仅 389 个，占比 1.5\%，而全球已注释的 P450 家族成员远超此数\cite{Nelson2009}，绝大多数尚无可用于机器学习训练的定量底物特异性数据。在如此不对称的数据分布下，通用预训练模型在 P450 家族上的 zero-shot 泛化能力成为一个尚未回答的关键问题。

与此同时，针对公开 checkpoint 的任何外部评估都面临数据泄漏的潜在风险：若测试集中的酶与原始训练集酶有重合，评估将无法区分真正的泛化能力与对已见样本的记忆效应。因此，针对 P450 家族建立泄漏控制的公平评估协议，并探究领域专属训练能否突破通用模型的泛化瓶颈，构成本论文的核心研究问题。

\section{研究目标实验设计与预期结果}
\label{sec:goals}

针对上述问题，本研究设想在 P450 超家族上对 EZSpecificity 模型进行系统的泛化能力测试，并在领域专属数据上从头训练同架构模型，以评估领域微调的必要性与收益；在此基础上进一步探索方向信息感知的架构改进。

为此本研究采用的方法链条为：（1）数据集构建：从 RCSB Protein Data Bank\cite{Burley2025,Berman2000}、ESIBank\cite{EZSpec2025}、P450Rdb、PlantP450DB\cite{Nelson2004} 与 PCPD\cite{Wang2021PCPD} 五个互补的数据源整合 P450 酶、底物与催化关系记录，经 UniProt\cite{UniProt2025} 序列桥接与 InChIKey\cite{InChI2015} 化学标识符去重合并，构建一个可复用的 P450 专属基准数据集；（2）模型训练与公平评估：在该基准上以与 EZSpecificity 同架构的模型从头训练基线，并以 InterPro\cite{Blum2025} 严格识别训练集重合酶、构造过滤后的外部测试集，与 EZSpecificity 公开 checkpoint 在泄漏控制下进行公平对比；（3）架构改进探索：在原架构基础上引入残基级几何向量感知（Geometric Vector Perceptron，GVP）通路，构造双图架构，探究方向信息编码对 P450 特异性预测的可行性。

研究所涉及的实验范围严格限于 P450 超家族的二分类底物特异性预测，采用单一随机切分第一折（random split fold~0）进行训练与评估，不开展多折交叉验证或多 seed 稳健性重复，相关限制在第六章局限性中如实说明。理论依据上，SOTA 模型在稀缺家族上的泛化不稳定是深度学习可重复性研究中的普遍发现；领域专属数据的微调可突破通用模型的分布偏差是已建立的方法学共识。预期结果为：EZSpecificity 公开模型在本研究外部测试集上将展现出明显的 zero-shot 泛化瓶颈；本研究构建的领域专属基线模型将显著突破该瓶颈；双图架构在单折评估下将展现对方向信息感知设计的初步可行性。

\section{论文组织}
\label{sec:arrangement}

本论文共分六章。第一章（本章）介绍研究背景、国内外研究现状、P450 家族上的特殊挑战与本研究的目标与实验设计。第二章综述酶底物特异性预测的相关工作，概述 EZSpecificity 的参考模型框架与所涉及的蛋白、分子表征学习组件，并明确本研究的问题定义与总体方案。第三章详述 P450 专属基准数据集的构建流程，包括五源数据获取、跨源去重对齐、双向负样本构造、数据切分、分子对接与特征生成。第四章报告原架构模型在 P450 专属数据上的训练结果，以及与 EZSpecificity 公开 checkpoint 的泄漏控制公平对比，并讨论其药学应用意义。第五章介绍双图架构改进的设计动机、残基级 GVP 模块的实现细节、单折训练与测试结果，以及相应的方法学启示。第六章总结本文的主要结论与创新贡献，如实声明研究局限性，并展望后续的方法学与数据层面改进方向。
